% !TeX root = bericht.tex
\section{Aufgabe 3}
\subsection{Aufgabenstellung}
In dieser Aufgabe soll eine Ofenregelung mit Matlab/Simulink modelliert und simuliert werden.
Die Regelstrecke sieht wie folgt aus:
\begin{figure}[H]
    \centering
    \includegraphics[width=\textwidth]{lab1_3_0.png}
    \caption{Zu bauende Ofenregelstrecke}
\end{figure}

Hierfür wurden in einer .m Datei die folgenden Parameter für unsere Strecke definiert:
\begin{sourcecode}[H]
    \caption{Parameter für die Übertragungsfunktion}
    \begin{minted}{matlab}
K1 = 10 / 400;      % 10V/400°C
K2 = 400 / 230;     % 400°C/230V
tau = 10 * 60;  
u_ein = 0.75;       
u_aus = -0.75;
    \end{minted}
\end{sourcecode}

\subsection{Modellierung in SimuLink}
\begin{figure}[H]
    \centering
    \includegraphics[width=\textwidth]{lab1_3_1.png}
    \caption{Simulink-Modell der Ofenregelstrecke}
\end{figure}
\subsubsection{Einstellung der Führungsgröße}
Der Ofen soll laut Aufgabenstellung auf 180°C für 30min aufgeheizt werden. 
Danach soll die Führungsgröße auf wieder auf 0°C gesetzt werden. 
Hierfür wurde in Simulink zwei Step-Blöcke verwendet (siehe Abbildung \ref{fig:steps}).\\
\begin{figure}[H]
    \centering
    \includegraphics[width=\textwidth]{scope_in.pdf}
    \caption{Einstellung der Führungsgröße mit Step-Blöcken}
    \label{fig:steps}
\end{figure} 
Durch die Addition der beiden Blöcke entsteht unsere gewünschte Führungsgröße 
mit einem Sprung von 0°C auf 180°C bei t=0s und einem Sprung von 180°C auf 0°C bei t=1800s (siehe Abbildung \ref{fig:fuehrung}).
Um die Führungsgröße in eine für den Regler brauchbare Spannung umzuwandeln, wurde ein Gain-Block mit einem Verstärkungsfaktor von $K_1$ verwendet.
\begin{figure}[H]
    \centering
    \includegraphics[width=\textwidth]{scope_f.pdf}
    \caption{Führungsgröße über die Zeit}
    \label{fig:fuehrung}
\end{figure}


\subsubsection{Modellierung des Reglers}
Der Regler selbst ist ein einfacher Relais, welche 230V oder 0V ausgibt, je nach Eingangssignal.
Das Eingangssignal wird aus der Differenz aus Führungsgröße und Istwert (über den Sensor gemessen) gebildet.
Der Relais-Block in Simulink hat die folgenden Einstellungen:
\begin{itemize}
    \item \bf\texttt{Switch on point}: \normalfont\texttt{u\_ein}
    \item \bf\texttt{Switch off point}: \normalfont\texttt{u\_aus}
    \item \bf\texttt{Output when on}: \normalfont\texttt{230}
    \item \bf\texttt{Output when off}: \normalfont\texttt{0}
\end{itemize}

\subsubsection{Modellierung der Heizung}
Die Heizung ist eine Übertragungsfunktion 
\begin{equation}
    G(s) = \frac{K_2}{1 + \tau s}
\end{equation}
Diese Übertragungsfunktion wurde mit dem gleichen Syntax wie in Quellcode \ref{code:tran_fun} beschrieben.

\subsubsection{Rückkopplung}
Die Rückkopplung der Strecke erfolgt über einen Sensor, welcher die aktuelle Temperatur misst und in eine Spannung umwandelt.
Hierfür wurde ein Gain-Block mit einem Verstärkungsfaktor von $K_1$ verwendet.

\subsection{Simulationsergebnisse}
Die Simulationsergebnisse für die Ofenregelstrecke sind in Abbildung \ref{fig:ofen} dargestellt.
Die gelbe Kurve zeigt die Führungsgröße, die orange Kurve zeigt die aktuelle Temperatur des Ofens und die blaue Kurve zeigt die Spannung, welche der Regler ausgibt.
\begin{figure}[H]
    \centering
    \includegraphics[width=\textwidth]{scope.pdf}
    \caption{Simulationsergebnisse der Ofenregelstrecke}
    \label{fig:ofen}
\end{figure}
Man kann hier klar erkennen, dass der Regler nicht in der Lage ist, die Führungsgröße von 180°C zu halten.
Die Temperatur steigt zwar an, aber da der Regler nur 230V oder 0V ausgeben kann, d.h. nur ein Ein/Ausschalten kann, 
kann die Temperatur nicht stabil gehalten werden und schwankt.
Da wir die An- und Ausschaltbedingungen des Relais auf $\pm 0.75V$ gesetzt haben, schaltet der Regler bei 210°C ein und bei 150°C aus.

\subsection{Variation von Ein- und Ausschalt-Threshold}
Die Schaltpunkte des Relais können durch die Parameter $u_{ein}$ und $u_{aus}$ angepasst werden.
Wenn wir $u_{ein}$/$u_{aus}$ niedriger machen, schaltet der Regler früher ein oder aus, was zu einer höheren Frequenz des Ein/Aus-Schaltens führt, also zu mehr kleineren Schwingungen).
Dies kann man klar in der folgenden Abbildung erkennen, in welcher $u_{ein}$ auf 0.25V und $u_{aus}$ auf -0.25V gesetzt wurden.
Der Regler sollte hier bei 190°C einschalten und bei 170°C ausschalten.
\begin{figure}[H]
    \centering
    \includegraphics[width=\textwidth]{scope_var.pdf}
    \caption{Simulationsergebnisse mit angepassten Schaltpunkten}
\end{figure}

Dieser Regler ist natürlich nicht optimal, da so eine hohe Frequenz den Relais beschädigen würde. 
Jedoch kann man es leicht als Einführung in SimuLink nehmen.
